% !TeX root = ../paper.tex
%
% Deutsche Zusammenfassung -- German-language abstract.
%
% REQUIRED: the supervisor confirmed by email (Sep 2026) that a German abstract is
% needed IN ADDITION to the English one, and that it does NOT count towards the
% page limit. It therefore sits on its own unnumbered page, before the title block,
% and page numbering starts at 1 on the body page that follows.
%
% This is a faithful translation of the English abstract in paper.tex, not a
% separate text. IF THE ENGLISH ABSTRACT CHANGES, CHANGE THIS TOO -- the two must
% not drift apart. Metric names are kept in their established forms
% ("Maintainability Index" is used untranslated in German SE writing; cyclomatic
% and cognitive complexity have standard German equivalents).

\thispagestyle{empty}

\section*{Zusammenfassung}

\noindent
Große Sprachmodelle erzeugen Python-Code, der häufig funktional korrekt ist.
Funktionale Korrektheit ist jedoch nur eine Teileigenschaft von Softwarequalität
und sagt wenig darüber aus, ob der erzeugte Code verständlich und wartbar ist. Die
vorliegende Arbeit untersucht, ob iteratives Feedback aus statischer Analyse die
Wartbarkeit von LLM-generiertem Python-Code verbessern kann, ohne dessen
funktionale Korrektheit zu verringern. Mit Qwen2.5-Coder-7B-Instruct, lokal unter
einem eingefrorenen, reproduzierbaren Protokoll ausgeführt (Greedy Decoding,
fester Seed), wurden 63 stratifizierte Aufgaben des MBPP+-Benchmarks generiert und
anschließend über bis zu drei Feedback-Iterationen verfeinert, in denen
strukturelle Messwerte (Maintainability Index, zyklomatische und kognitive
Komplexität) als erneute Prompts an das Modell zurückgegeben wurden. Verglichen
wurden drei Konfigurationen der Schleife: zwei Varianten ohne korrektheitsbasiertes
Zurücksetzen, die sich allein darin unterscheiden, ob beim Refactoring Kommentare
entfernt werden dürfen, sowie eine Variante, die jedes Refactoring
verwirft und zurücksetzt, das das MBPP+-Orakel verletzt. Die kognitive Komplexität
sank in jeder Konfiguration signifikant, und dieses Ergebnis hält einer
Bonferroni-Korrektur über die zwölf primären Tests (drei Konfigurationen mal vier
Zielgrößen) stand. Die funktionale
Korrektheit zeigte dabei keinen statistisch signifikanten Verlust; die zurücksetzende Variante schließt
Korrektheitsregressionen konstruktionsbedingt aus. Der
Maintainability Index sank in ähnlichem Maße, unabhängig davon, ob Kommentare
entfernt oder erhalten wurden. Dieser Rückgang ist auf seinen nicht-monotonen
Kommentarterm zurückzuführen und nicht auf einen Verlust struktureller Qualität,
was den Index hier als Ergebnismetrik ungeeignet macht. Die Ergebnisse
zeigen, dass Feedback aus statischer Analyse strukturelle Komplexität messbar
reduzieren kann und dass eine schlanke Korrektheitsprüfung das Verfahren in diesem
Rahmen sicher anwendbar macht.

\vspace{0.8em}
\noindent\textbf{Schlagwörter:} LLM-Codegenerierung, Software-Wartbarkeit,
kognitive Komplexität, statische Analyse, Feedback-Schleife, MBPP
